\documentclass[10pt,a4paper]{article}

\usepackage{fontspec}
\usepackage{polyglossia}
\setmainlanguage{russian}
\setotherlanguage{english}
\setmainfont{DejaVu Serif}
\setsansfont{DejaVu Sans}
\setmonofont{DejaVu Sans Mono}
\usepackage{graphicx}
\usepackage{geometry}
\geometry{a4paper, margin=2cm}

\usepackage{amsmath,amssymb,amsfonts}
\usepackage{enumitem}
\usepackage[table]{xcolor}
\usepackage{tcolorbox}
\tcbuselibrary{breakable, skins}
\usepackage{titlesec}
\usepackage{fancyhdr}
\usepackage{parskip}
\usepackage{tocloft}

\renewcommand{\cftsecleader}{\cftdotfill{\cftdotsep}}
\renewcommand{\cftsubsecleader}{\cftdotfill{\cftdotsep}}
\renewcommand{\cftsubsubsecleader}{\cftdotfill{\cftdotsep}}

\renewcommand{\cftsecpagefont}{\normalfont}
\renewcommand{\cftsubsecpagefont}{\normalfont}
\renewcommand{\cftsubsubsecpagefont}{\normalfont}

\setlength{\cftsecnumwidth}{2.5em}
\setlength{\cftsubsecnumwidth}{3.2em}
\setlength{\cftsubsubsecnumwidth}{4.2em}

\setlength{\cftsecindent}{0em}
\setlength{\cftsubsecindent}{1.5em}
\setlength{\cftsubsubsecindent}{3em}

\usepackage[colorlinks=true, linkcolor=blue!50!black, citecolor=blue!50!black, urlcolor=blue!50!black]{hyperref}

\definecolor{yesgreen}{RGB}{20,110,60}
\definecolor{nored}{RGB}{160,30,30}
\definecolor{sectionblue}{RGB}{20,40,90}
\definecolor{boxgray}{RGB}{245,246,248}
\definecolor{ruleblue}{RGB}{60,90,150}

\definecolor{mainblue}{RGB}{30,60,110}
\definecolor{boxbg}{RGB}{237,242,250}
\definecolor{boxframe}{RGB}{140,170,210}
\definecolor{intuitbg}{RGB}{255,247,230}
\definecolor{intuitframe}{RGB}{220,180,90}

\titleformat{\section}
  {\normalfont\Large\bfseries\color{sectionblue}}
  {\thesection}{0.6em}{}[{\color{ruleblue}\titlerule[1pt]}]
\titleformat{\subsection}
  {\normalfont\large\bfseries\color{sectionblue}}
  {\thesubsection}{0.6em}{}

\pagestyle{fancy}
\fancyhf{}
\fancyhead[L]{\small\textit{Макроэкономика финансовых рынков. Банк тестов}}
\fancyhead[R]{\small\textit{ИП ФЭН, 2026-2027}}
\fancyfoot[C]{\small\thepage}
\renewcommand{\headrulewidth}{0.4pt}

\newcommand{\qda}{\textbf{\color{yesgreen}Ответ: ДА}}
\newcommand{\qnet}{\textbf{\color{nored}Ответ: НЕТ}}

\newtcolorbox{qbox}[1]{
  breakable, enhanced, colback=boxgray, colframe=ruleblue!70!black,
  boxrule=0.5pt, left=7pt, right=7pt, top=6pt, bottom=6pt,
  arc=3pt, title={#1}, fonttitle=\bfseries, coltitle=black,
  colbacktitle=ruleblue!12!white, attach boxed title to top left,
  boxed title style={arc=2pt, boxrule=0pt}
}

\newenvironment{answer}[2]{%
  \begin{qbox}{Утверждение #1}
  \textit{#2}\par\smallskip
}{%
  \end{qbox}\vspace{2pt}
}

\title{\color{sectionblue}\Large\textbf{Макроэкономика финансовых рынков}\\[4pt]
\large Разбор банка вопросов}
\date{}

\begin{document}
\begin{titlepage}
    \thispagestyle{empty}
    \centering

    \textcolor{mainblue}{\rule{\textwidth}{1.2pt}}

    \vspace{1.2cm}

    {\color{mainblue}\Huge\bfseries Макроэкономика\\[4pt] финансовых рынков\par}

    \vspace{0.8cm}

    {\color{mainblue}\Large Модуль 1\\[2pt] Динамические модели общего равновесия\par}

    \vspace{0.6cm}

    \textcolor{mainblue}{\rule{\textwidth}{0.4pt}}

    \vfill

    \begin{minipage}{0.85\textwidth}
        \small\itshape
        Данный разбор банка тестовых вопроосов для квизов по макро-2 подготовлен студенткой ОП "Экономика" Александрой Подлеснюк
    \end{minipage}

    \vspace{1cm}

    \textcolor{mainblue}{\rule{\textwidth}{1.2pt}}
\end{titlepage}

\newpage
\tableofcontents
\newpage

\section{Часть 1. Межвременная оптимизация потребительского выбора}

\begin{answer}{1}{Динамическое бюджетное ограничение домохозяйства описывает динамику его активов или долга.}
\qda. Динамическое бюджетное ограничение задаётся уравнением накопления чистых финансовых активов:
\[
A_{t+1} = (A_t + Y_t - C_t)(1+r) = (A_t + S_t)(1+r)
\]
Оно показывает, как накопленное к началу периода t богатство $A_t$, вместе с несбереженным доходом текущего периода $(S_t = Y_t − C_t)$, реинвестированное по ставке r, превращается в богатство следующего периода $A_{t+1}$. Если $A_t$ < 0, переменная интерпретируется как долг.В отличие от статического ограничения, связывающего расходы с доходом в один момент времени, данное уравнение описывает трансформирование запаса богатства во времени через накопление сбережений и начисление процентов. Таким образом, ДБО - это именно закон движения запаса активов (или долга) во времени, а не условие, определяющее сам уровень потребления
\end{answer}

\begin{answer}{2}{Динамическое бюджетное ограничение домохозяйства — это уравнение динамики его потребления.}
\qnet. Уравнение $A_{t+1} = (A_t + Y_t - C_t)(1+r)$ является уравнением движения переменной состояния — активов $A_t$. Динамическое бюджетное ограничение
домохозяйства описывает динамику его активов или долга, показывает изменение
богатства путем накопления сбережений и получения дохода (могут быть
отрицательными). Потребление $C_t$ выступает в нём как управляющая переменная. Динамика самого потребления во времени определяется не бюджетным ограничением, а условием первого порядка задачи максимизации полезности - правилом Рамсея–Кейнса (уравнением Эйлера)
\end{answer}

\begin{answer}{3}{Динамическое бюджетное ограничение домохозяйства связывает его потребление и доход с приростом в сбережениях.}
\qnet. Уравнение $A_{t+1} = (A_t + S_t)(1+r)$ связывает изменение финансового богатства/долга между периодами $t$ и $t+1$ с текущим запасом активов $A_t$ и объёмом сбережений $S_t = Y_t - C_t$, дисконтированных по процентной ставке $r$. Доход $Y_t$ в этом уравнении является экзогенным или эндогенно определяемым потоком, но не величиной, прирост которой определяется бюджетом
\end{answer}

\begin{answer}{4}{Динамическое бюджетное ограничение домохозяйства связывает его потребление и богатство/долг с приростом в доходе.}
\qnet. Динамическое бюджетное ограничение связывает изменение активов/долга домохозяйства в двух периодах в реальном выражении (в зависимости от ставки процента), в зависимости от богатства и сбережений в периоде $t$. Оно не описывает динамику или прирост дохода $Y_t$, который задаётся экзогенно здесь
\end{answer}

\begin{answer}{5}{Условие отсутствия игр Понци требует, чтобы приведённая стоимость богатства агента в последнем периоде жизни (при конечном горизонте) равнялась нулю.}
\qda. Для конечного временного горизонта $T$ (с последним периодом $T-1$) условие No-Ponzi-Game (NPG) требует $A_T = 0$. В приведённых к текущему периоду $t$ ценах это выражается как:
\[
\frac{A_T}{(1+r)^{T-t}} = 0
\]
Рациональному потребителю не стоит оставлять каких-либо активов после своей
смерти, предполагая, что нет мотива оставлять наследство, альтруизма (конечный
временной горизонт). Рациональным кредиторам не следует давать в долг потребителю,
который не сможет выплатить долг до конца жизни
\end{answer}

\begin{answer}{6}{Рациональный агент не должен оставлять богатство после смерти (в отсутствие мотива оставления наследства и прочих соображений).}
\qda. Если после смерти агента остаётся положительное богатство $A_T > 0$, такое поведение нерационально: перенаправив эти ресурсы на потребление в течение жизни, агент получил бы строго большую суммарную полезность, так как $u'(C) > 0$
\end{answer}

\begin{answer}{7}{Рациональный инвестор не должен давать в долг агенту, если тот не расплатится по долгам до смерти.}
\qda. Наличие непогашенного долга $A_T < 0$ означает убыток для кредитора. Рациональный кредитор, стремящийся к максимизации собственной полезности и возврату заёмных средств с процентами, откажет в предоставлении кредита агенту, не способному расплатиться до конца жизни
\end{answer}

\begin{answer}{8}{Межвременное бюджетное ограничение домохозяйства требует: в каждый момент времени приведённая стоимость расходов на потребление равна приведённой стоимости будущих доходов (с учётом налогов и накопленного богатства/долга).}
\qda. Последовательно итерируя динамическое ограничение $A_{\tau+1} = (A_\tau + Y_\tau^d - C_\tau)(1+r)$ от $\tau = t$ до $T-1$ и накладывая условие $A_T = 0$, получаем межвременное бюджетное ограничение:
\[
\sum_{\tau=t}^{T-1} \frac{C_\tau}{(1+r)^{\tau-t+1}} = A_t + \sum_{\tau=t}^{T-1} \frac{Y_\tau^d}{(1+r)^{\tau-t+1}}
\]
где $Y_\tau^d = Y_\tau - T_\tau$ — располагаемый доход. Текущие активы $A_t$ плюс дисконтированный поток доходов за вычетом налогов должны строго равняться дисконтированному потоку расходов на потребление. 
\textbf{Интуиция:} предположим, что это не так и потребление меньше, тогда это неоптимальное поведение, поскольку доступное увеличение потребления вызвало бы рост полезности. А если предположить, что потребление больше, чем дисконтированая стоимость будущих доходов, то у потребителя нет ресурсов на это, а заем невозможен, так как нарушается условие отсутствия игр Понци
\end{answer}

\begin{answer}{9}{Межвременное бюджетное ограничение домохозяйства требует: в каждый момент времени текущие расходы на потребление должны быть равны текущему доходу (с учётом налогов и накопленного богатства/долга).}
\qnet. Межвременное ограничение связывает суммарные дисконтированные потоки за весь жизненный цикл. Домохозяйство может перераспределять ресурсы во времени через финансовые рынки, поэтому в отдельные периоды $C_t > Y_t^d$ (за счёт займов или растраты накопившихся активов) или $C_t < Y_t^d$ (за счёт сбережений или возможности дать средства в долг) 
\end{answer}

\begin{answer}{10}{Если в текущем периоде домохозяйство имеет положительное богатство, оно может позволить себе отрицательную приведённую стоимость сбережений в будущем.}
\qda. Из межвременного ограничения следует: $A_t = -\sum_{\tau=t}^{T-1} \frac{S_\tau}{(1+r)^{\tau-t+1}}$, где $S_\tau = Y_\tau^d - C_\tau$. Если $A_t > 0$, то сумма будущих дисконтированных сбережений должна быть отрицательной ($\sum \frac{S_\tau}{(1+r)^{\tau-t+1}} < 0$). Это означает, что агент будет потреблять больше текущего дохода, «проедая» накопленное богатство
\end{answer}

\begin{answer}{11}{Если в текущем периоде домохозяйство является должником, обеспечением его долга должна выступать положительная приведённая стоимость будущих сбережений.}
\qda. Если $A_t < 0$, то для выполнения межвременного ограничения требуется $\sum_{\tau=t}^{T-1} \frac{S_\tau}{(1+r)^{\tau-t+1}} = -A_t > 0$. Домохозяйство обязано в будущем сберегать в среднем больше, чем потреблять ($Y^d > C$), чтобы погасить текущий долг и накопившиеся проценты
\end{answer}

\begin{answer}{12}{Если в текущем периоде домохозяйство имеет положительное богатство, обеспечением его богатства должна выступать положительная приведённая стоимость будущих сбережений.}
\qnet. Наоборот, при $A_t > 0$ само богатство выступает источником финансирования будущего потребления, поэтому суммарная приведённая стоимость будущих сбережений должна быть отрицательной
\end{answer}

\begin{answer}{13}{Если в текущем периоде домохозяйство является должником, обеспечением его долга должна выступать отрицательная приведённая стоимость будущих сбережений.}
\qnet. Если при $A_t < 0$ будущие сбережения будут отрицательными ($C > Y^d$), долг агента будет экспоненциально расти, что нарушит условие NPG  он будет потреблять больше, чем ему позволяют активы и доход, à приведенная стоимость будущих сбережений не может быть отрицательной у рационального агента, имеющего долг перед рациональным кредитором)
\end{answer}

\begin{answer}{14}{Условие отсутствия игр Понци в модели с бесконечным временным горизонтом требует: приведённая стоимость богатства (долга) домохозяйства должна стремиться к нулю на бесконечном временном горизонте.}
\qda. В модели с бесконечным горизонтом условие NPG имеет вид:
\[
\lim_{T \to \infty} \frac{A_T}{(1+r)^T} = 0
\]
Если предел отрицателен, агент бесконечно наращивает долг, что недопустимо для кредиторов (так как это значит, что кредиторы были
нерациональными, раз выдали деньги ненадежному заемщику). Если предел положителен, агент недопотребляет, что нарушает оптимизацию полезности и  несовместимо с предпосылкой о рациональности агентов
\end{answer}

\begin{answer}{15}{Условие отсутствия игр Понци в модели с бесконечным временным горизонтом требует: текущая стоимость богатства (долга) домохозяйства должна стремиться к нулю на бесконечном временном горизонте.}
\qnet. Требование накладывается на дисконтированную стоимость $\frac{A_T}{(1+r)^T}$, а не на номинальный объем богатства $A_T$. Величина $A_T$ может возрастать во времени (например, линейно), при этом приведённая стоимость будет стремиться к нулю. Поэтому необходимо рассматривать именно приведенную стоимость, так как рассматриваем вопрос о
целесообразности накопления богатства с точки зрения текущего периода
\end{answer}

\begin{answer}{16}{Условие отсутствия игр Понци в модели с бесконечным временным горизонтом требует: богатство не должно расти «слишком быстро» — темп прироста богатства (долга) не должен систематически превышать ставку процента.}
\qda. Пусть $A_T = A_0 (1+g_A)^T$. Предел приведённой стоимости равен:
\[
\lim_{T \to \infty} \frac{A_0 (1+g_A)^T}{(1+r)^T} = A_0 \lim_{T \to \infty} \left(\frac{1+g_A}{1+r}\right)^T = 0 \iff g_A < r
\]
Фактически, если темп роста долга $g_A \ge r$, то заёмщик может вечно финансировать выплаты по старым долгам за счёт эмиссии новых, создавая финансовую пирамиду, что запрещено условием NPG. Если темп роста богатства будет выше ставки, это будет означать, что целесообразно накапливать богатство с точки зрения будущего периода, то есть предел не будет стремиться к нулю
\end{answer}

\begin{answer}{17}{Условие отсутствия игр Понци в модели с бесконечным временным горизонтом требует: если у агента есть долг, он рано или поздно должен быть полностью погашен.}
\qnet. Условие NPG не требует полного обнуления номинального долга ($A_T = 0$) в какой-либо фиксированный момент времени. Долг может оставаться отрицательным и даже неограниченного расти по абсоллютному значению - это не нарушает NPG, важно лишь, чтобы темп прироста величины долга не превышал ставку процента
\end{answer}

\begin{answer}{18}{Правило Рамсея-Кейнса — это условие первого порядка в задаче межвременной оптимизации поведения домохозяйства.}
\qda. В задаче максимизации аддитивно-сепарабельной полезности:
\[
\max_{\{C_t\}} \sum_{t=0}^{T-1} \frac{u(C_t)}{(1+\rho)^t} \quad \text{при} \quad A_{t+1} = (A_t + Y_t^d - C_t)(1+r)
\]
записывая функции Лагранжа и беря частные производные по $C_t$ и $A_{t+1}$, получаем условие первого порядка (уравнение Эйлера):
\[
\frac{u'(C_t)}{u'(C_{t+1})} = \frac{1+r}{1+\rho}
\]
\end{answer}

\begin{answer}{19}{Правило Рамсея-Кейнса — это условие оптимального межвременного выбора потребительских расходов.}
\qda. Данное правило задаёт оптимальную динамику потребления во времени, связывая отношение предельных полезностей в соседних периодах со нормой доходности $r$ и нормой субъективных межвременных предпочтений $\rho$
\end{answer}

\begin{answer}{20}{Оптимальный межвременной выбор потребления требует равенства полезностей от текущего и будущего потребления.}
\qnet. Оптимум требует равенства отношения предельных полезностей величине $(1+r)/(1+\rho)$, а не равенства самих (уровней) полезностей текущего и будущего потребления. Равенство предельных полезностей $u'(C_t) = u'(C_{t+1})$ и самих уровней потребления достигается только в частном случае $r = \rho$
\end{answer}

\begin{answer}{21}{Оптимальный межвременной выбор потребления требует: с учётом субъективного дисконтирования, отношение предельных полезностей потребления в двух близлежащих периодах в оптимуме должно равняться отношению соответствующих стоимостей.}
\qda. Нормируя цену потребления в периоде $t$ к $P_t = 1$, относительная стоимость будущего потребления составляет $P_{t+1} = \frac{1}{1+r}$. Отношение дисконтированных предельных полезностей (MRS) равно:
\[
\text{MRS}_{t, t+1} = \frac{u'(C_t)}{u'(C_{t+1})/(1+\rho)} = \frac{P_t}{P_{t+1}} = 1+r \implies \frac{u'(C_t)}{u'(C_{t+1})} = \frac{1+r}{1+\rho}
\]
Это классическое условие равенства предельной нормы замещения соотношению цен (отношение предельных норм замещения равно отношению цен потребления в соседних периодах)
\end{answer}

\begin{answer}{22}{Если в задаче межвременного выбора потребления нормировать стоимость единицы потребления в первом периоде к 1, тогда $\frac{1}{1+r}$ является стоимостью единицы потребления во втором периоде.}
\qda. Из межвременного бюджетного ограничения $C_1 + \frac{C_2}{1+r} = Y_1^d + \frac{Y_2^d}{1+r}$ следует, что при нормировке цены $C_1$ к 1, относительная стоимость единицы потребления во втором периоде выраженная в единицах первого периода равна $\frac{1}{1+r}$
\end{answer}

\begin{answer}{23}{Увеличение субъективной нормы межвременных предпочтений снизит отношение предельных полезностей потребления в первом и во втором периоде, приводя к увеличению потребления в первом периоде и снижению во втором.}
\qda. Рост $\rho$ снижает правую часть уравнения Эйлера $\frac{1+r}{1+\rho}$. Для восстановления равенства отношение $\frac{u'(C_t)}{u'(C_{t+1})}$ должно уменьшиться. При убывающей предельной полезности это означает рост $C_t$ и снижение $C_{t+1}$ - домохозяйство становится «нетерпеливее» и смещает потребление к текущему периоду
\end{answer}

\begin{answer}{24}{Увеличение субъективной нормы межвременных предпочтений увеличит отношение предельных полезностей потребления в первом и во втором периоде, приводя к снижению потребления в первом периоде и увеличению во втором.}
\qnet. Рост $\rho$ снижает правую часть $\frac{1+r}{1+\rho}$, что ведёт к росту текущего потребления за счёт будущего, а не наоборот
\end{answer}

\begin{answer}{25}{Увеличение ставки процента снижает относительную стоимость потребления во втором периоде и увеличивает отношение предельной полезности потребления в первом периоде к предельной полезности потребления во втором периоде.}
\qda. Относительная цена будущего потребления равна $\frac{1}{1+r}$, она отрицательно зависит от $r$. В уравнении Эйлера $\frac{u'(C_t)}{u'(C_{t+1})} = \frac{1+r}{1+\rho}$ рост $r$ увеличивает правую часть и, следовательно, увеличивает отношение предельных полезностей.
\end{answer}

\begin{answer}{26}{Увеличение ставки процента увеличивает относительную стоимость потребления во втором периоде и снижает отношение предельной полезности потребления в первом периоде к предельной полезности потребления во втором периоде.}
\qnet. Рост $r$ снижает цену будущего потребления $\frac{1}{1+r}$ и увеличивает соотношение предельных полезностей $\frac{u'(C_t)}{u'(C_{t+1})}$
\end{answer}

\begin{answer}{27}{С учётом убывания предельной полезности, увеличение ставки процента приведёт к желанию потреблять меньше в первом периоде и больше во втором.}
\qda. Повышение $r$ увеличивает отношение $\frac{u'(C_t)}{u'(C_{t+1})} = \frac{1+r}{1+\rho}$. При $u'' < 0$ рост предельной полезности $u'(C_t)$ требует сокращения текущего потребления $C_t \downarrow$, а падение $u'(C_{t+1})$ — увеличения будущего потребления $C_{t+1} \uparrow$
\end{answer}

\begin{answer}{28}{С учётом убывания предельной полезности, увеличение ставки процента приведёт к желанию потреблять больше в первом периоде и меньше во втором.}
\qnet. Рост $r$ повышает доходность сбережений, удорожая текущее потребление и стимулируя его снижение в пользу будущего
\end{answer}

\begin{answer}{29}{Увеличение ставки процента увеличивает капитальный доход во втором периоде и, следовательно, общий доход на протяжении жизни.}
\qda. При наличии сбережений в первом периоде ($S_t > 0$) доход от процентов во втором периоде составляет $r S_t$. Повышение $r$ увеличивает этот капитальный доход в периоде $t+1$, расширяя общее межвременное бюджетное ограничение (эффект дохода)
\end{answer}

\begin{answer}{30}{Эффекты межвременного замещения и дохода характеризуются противоположным воздействием на потребление в первом периоде.}
\qda. При росте $r$ для кредитора ($S_t > 0$):
\begin{itemize}[leftmargin=1.5em]
  \item \textbf{Эффект замещения (SE):} текущее потребление относительно удорожается, что побуждает сокращать $C_t \downarrow$ и увеличивать сбережения $S_t \uparrow$ и  потребление в следующем периоде $C_{t+1} \uparrow$
  \item \textbf{Эффект дохода (IE):} рост процентного дохода $r S_t$ увеличивает совокупное богатство, стимулируя рост потребления во всех периодах, то есть $C_t \uparrow$, $C_{t+1} \uparrow$ и при этм $S_t \downarrow$
\end{itemize}
Таким образом, SE и IE действуют на $C_t$ в противоположных направлениях
\end{answer}

\begin{answer}{31}{Эффекты межвременного замещения и дохода характеризуются противоположным воздействием на сбережения в первом периоде.}
\qda. Поскольку $S_t = Y_t^d - C_t$, эффект замещения снижает $C_t$ (увеличивает $S_t \uparrow$), а эффект дохода увеличивает $C_t$ (снижает $S_t \downarrow$). Воздействие на сбережения строго противоположно
\end{answer}

\begin{answer}{32}{Эффекты межвременного замещения и дохода характеризуются противоположным воздействием на потребление во втором периоде.}
\qnet. Для будущего потребления $C_{t+1}$ оба эффекта действуют сонаправленно: SE стимулирует рост $C_{t+1} \uparrow$ (так как общее будущее потребление подешевело), и IE также увеличивает $C_{t+1} \uparrow$ из-за выросшего богатства
\end{answer}

\begin{answer}{33}{Модель межвременного выбора предполагает, что эффект межвременного замещения доминирует над эффектом дохода.}
\qda. В стандартных макроэкономических предпосылках эластичность межвременного замещения достаточно велика ($\sigma > 1$, $\sigma = \frac{u '(c)}{u''(c)c} = \lim_{\tau \to t} \sigma(c_t, c_{\tau})$), благодаря чему эффект замещения доминирует над эффектом дохода, гарантируя положительный наклон функции сбережений по процентной ставке ($\frac{\partial S_t}{\partial r} > 0$)
\end{answer}

\begin{answer}{34}{Если $r > \rho$, потребление репрезентативного агента будет возрастать на протяжении жизненного цикла.}
\qda. Из уравнения Эйлера $\frac{u'(C_t)}{u'(C_{t+1})} = \frac{1+r}{1+\rho}$. При $r > \rho$ имеем $\frac{1+r}{1+\rho} > 1$, откуда $u'(C_t) > u'(C_{t+1})$. В силу строгой вогнутости функции полезности ($u'' < 0$) это означает $C_t < C_{t+1}$.И это верно для любой пары соседних периодов, то есть на протяжении всей жизни, поэтому потребление монотонно возрастает во времени
\end{answer}

\begin{answer}{35}{Если $r > \rho$, потребление репрезентативного агента будет снижаться на протяжении жизненного цикла.}
\qnet. При $r > \rho$ потребление монотонно растет, а не снижается
\end{answer}

\begin{answer}{36}{Если $r < \rho$, потребление репрезентативного агента будет снижаться на протяжении жизни.}
\qda. При $r < \rho$ получаем $\frac{1+r}{1+\rho} < 1 \implies u'(C_t) < u'(C_{t+1}) \implies C_t > C_{t+1}$. Высокая субъективная нетерпеливость приводит к монотонному падению потребления во времени
\end{answer}

\begin{answer}{37}{Если $r < \rho$, потребление репрезентативного агента будет возрастать на протяжении жизни.}
\qnet. При $r < \rho$ потребление убывает во времени
\end{answer}

\begin{answer}{38}{Если $r = \rho$, оптимальное поведение репрезентативного агента характеризуется выбором постоянного уровня потребления на протяжении всей жизни.}
\qda. При $r = \rho$ имеем $\frac{1+r}{1+\rho} = 1 \implies u'(C_t) = u'(C_{t+1})$. Из строгой монотонности $u'(\cdot)$ следует сглаживание потребления: $C_0 = C_1 = \dots = C_T$
\end{answer}

\begin{answer}{39}{В соответствии с теорией межвременного выбора тренд в потребительских расходах может быть объяснён различием между ставкой процента и нормой субъективных межвременных предпочтений, но не траекторией дохода.}
\qda. Наклон (тренд) траектории потребления задаётся исключительно соотношением $r$ и $\rho$ через уравнение Эйлера. Поток доходов $Y_t$ определяет лишь абсолютный уровень всей траектории потребления через межвременное бюджетное ограничение, но не её форму/наклон
\end{answer}

\begin{answer}{40}{В соответствии с теорией межвременного выбора тренд в потребительских расходах может быть объяснён траекторией дохода, но не различием между ставкой процента и нормой субъективных межвременных предпочтений.}
\qnet. Тренд определяется лишь соотношением между нормой доходности и нормой субъективных межвременных предпочтений. Траектория дохода влияет на уровень потребления (через межвременное бюджетное ограничение), но не на его тренд (форму траектории)
\end{answer}

\begin{answer}{41}{В соответствии с теорией межвременного выбора потребитель всегда выбирает постоянный уровень потребления.}
\qnet. Постоянное потребление выбирается исключительно в частном случае $r = \rho$. При $r \neq \rho$ оптимальная траектория потребления является монотонно возрастающей или убывающей
\end{answer}

\newpage
\section{Часть 2. Рикардианская эквивалентность, концепции Эйбела и Кругмана}

\begin{answer}{1}{В соответствии с рикардианской эквивалентностью, выбор между финансированием определённого потока государственных закупок за счёт налогообложения или государственных заимствований иррелевантен для определения потребления.}
\qda. Подстановка межвременного бюджетного ограничения правительства в ограничение домохозяйств приводит к консолидированному бюджетному ограничению экономики, в которое не входят ни государственные облигации $b_t$, ни траектория налоговых поступлений $T_{t+\tau}$.

\begin{enumerate}
    \item \textbf{Межвременное бюджетное ограничение домохозяйств:}
    Предположим, что государственные облигации являются единственным видом активов в экономике, то есть $A_t = b_t$. Межвременное бюджетное ограничение домохозяйства имеет вид:
    \[
    \sum_{\tau=0}^{\infty} \frac{C_{t+\tau}}{(1+r)^\tau} = A_t + \sum_{\tau=0}^{\infty} \frac{Y_{t+\tau} - T_{t+\tau}}{(1+r)^\tau} \quad 
    \]

    \item \textbf{Межвременное бюджетное ограничение государства:}
    Приведенная стоимость государственных закупок с учетом имеющегося государственного долга $b_t$ должна равняться приведенной стоимости налоговых поступлений:
    \[
    \sum_{\tau=0}^{\infty} \frac{G_{t+\tau}}{(1+r)^\tau} + b_t = \sum_{\tau=0}^{\infty} \frac{T_{t+\tau}}{(1+r)^\tau} \quad
    \]

    \item \textbf{Консолидация бюджетных ограничений:}
    Подставляя выражение для налогов в ограничение домохозяйств и учитывая, что $A_t = b_t$:
    \[
    \sum_{\tau=0}^{\infty} \frac{C_{t+\tau}}{(1+r)^\tau} = b_t + \sum_{\tau=0}^{\infty} \frac{Y_{t+\tau}}{(1+r)^\tau} - \underbrace{\left( b_t + \sum_{\tau=0}^{\infty} \frac{G_{t+\tau}}{(1+r)^\tau} \right)}_{\sum \frac{T_{t+\tau}}{(1+r)^\tau}}
    \]
    Раскрывая скобки, находим консолидированное бюджетное ограничение:
    \[
    \sum_{\tau=0}^{\infty} \frac{C_{t+\tau}}{(1+r)^\tau} = \sum_{\tau=0}^{\infty} \frac{Y_{t+\tau} - G_{t+\tau}}{(1+r)^\tau} \quad
    \]

    \item \textbf{Рикардианская эквивалентность:}
    Из итогового уравнения видно, что государственные облигации $b_t$ и активы $A_t$ полностью взаимосокращаются, а потоки налогов $T_{t+\tau}$ выпадают из уравнения. Приведенная стоимость потребления домохозяйств определяется исключительно приведенной базой доходов за вычетом государственных закупок $\{Y_{t+\tau} - G_{t+\tau}\}_{\tau=0}^{\infty}$. Следовательно, выбор между текущим налогообложением или заменой налогов государственными займами (облигациями) не меняет межвременного богатства домохозяйств и не оказывает влияния на траекторию потребления
\end{enumerate}
\end{answer}

\begin{answer}{2}{В соответствии с рикардианской эквивалентностью, ограничением на потребление выступает приведённая сумма будущих государственных закупок, а не способ их финансирования.}
\qda. Объединённое бюджетное ограничение содержит поток реальных государственных закупок $G_\tau$, а не отдельные величины налогов $T_\tau$ или долга $B_\tau$
И получается, что бюджетное множество домохозяйств зависит только от приведённой стоимости государственных закупок с учетом активов, а не от того, каким набором налогов или долгов эти закупки профинансированы
\end{answer}

\begin{answer}{3}{В соответствии с рикардианской эквивалентностью, фискальная политика в принципе не должна оказывать воздействие на потребление.}
\qnet. Эквивалентность утверждает нейтральность лишь способа финансирования данной траектории госзакупок. Само изменение уровня и траектории $G$ напрямую входит в бюджетное ограничение и меняет потребление домохозяйств. Также фискальная политика влияет, если налоги являются функцией от результата экономической деятельности, а не фиксированными единовременными изьятиями (тк такие налоги не искажающие).В результате чего меняется экономическая активность и поэтому меняется будущий доход
\end{answer}

\begin{answer}{4}{В соответствии с рикардианской эквивалентностью, домохозяйства не должны рассматривать владение государственными облигациями как накопление чистого богатства.}
\qda. Облигации $B_t$ приносят процентный доход, однако обеспечением государственного долга выступают будущие налоги, которые заплатят те же самые (впередсмотрящие, рациональные) домохозяйства. Приобретая сегодня государственную облигацию (финансируя снижение налогов), рациональный агент должен понимать, что в будущем ему (или его наследникам, при наличии альтруизма) придётся заплатить дополнительные налоги ровно на приведённую стоимость этого долга, поэтому чистое богатство (с учётом обязательств по будущим налогам) не меняется
\end{answer}

\begin{answer}{5}{Рикардианская эквивалентность возникает только в моделях, где домохозяйства имеют бесконечный временной горизонт. В реальности домохозяйства имеют конечный временной горизонт. И это неизбежно нарушает рикардианскую эквивалентность.}
\qnet. В базовой модели с бесконечным горизонтом планирования репрезентативного агента момент будущего повышения налогов действительно не имеет значения. При конечном же горизонте без межпоколенческих связей РЭ локально нарушается: если повышение налогов ожидается после смерти текущего поколения, агенты воспринимают временное снижение налогов как прирост своего постоянного дохода и увеличивают текущее потребление.

Однако сам по себе факт конечности жизни отдельного индивида не обязан нарушать нейтральность фискальной политики. Как показал Р. Барро (1974) в рамках модели перекрывающихся поколений (OLG), при наличии межпоколенческого альтруизма конечный жизненный цикл отдельного домохозяйства превращается в  бесконечный горизонт планирования династии.

Предположим, что полезность родителей напрямую зависит от полезности детей:
\begin{equation}
    U_t = u(C_t) + \beta U_{t+1}, \quad \beta \in (0, 1)
\end{equation}
где $\beta$ это коэффициент субъективного межпоколенческого дисконтирования. В этом случае родительское поколение рассматривает благосостояние потомков как часть собственного. При снижении налогов сегодня (финансируемом выпуском государственного долга) родители понимают, что неизбежное будущее увеличение налогового бремени падет на их детей. Чтобы сгладить межвременное потребление династии (с учетом субъективного дисконтирования $\beta$ и темпов роста населения), родители не увеличивают свое текущее потребление, а сберегают высвободившиеся средства и увеличивают размер передаваемого наследства на величину будущих налоговых обязательств потомков.

В результате текущее поколение не рассматривает будущие налоги как чужую проблему, и аргумент о конечности временного горизонта перестает опровергать принцип нейтральности фискальной политики.

Эквивалентность Рикардо сохраняется только при условии внутреннего решения оптимизационной задачи, когда мотив оставления наследства является оперативным (действующим), то есть оптимальный объем передаваемого наследства строго положителен ($B_t > 0$).

РЭ перестает работать и возникает угловое решение ($B_t = 0$) в следующих случаях:
\begin{enumerate}
    \item \textbf{Высокий темп роста населения:} если у родителей много детей, распределение ресурсов на каждого потомка существенно размывает эффект наследства.
    \item \textbf{Сильное дисконтирование потомков:} если родители недостаточно высоко оценивают полезность детей (малое значение $\beta$)
\end{enumerate}

В этих обстоятельствах родители предпочли бы оставить отрицательное наследство (то есть занять у будущих поколений), однако институциональное ограничение $B_t \ge 0$ не позволяет этого сделать. Ограничение становится связывающим ($B_t = 0$), родители оставляют нулевое наследство, и текущее снижение налогов полностью уходит в их собственное потребление, поэтому рикардианская эквивалентность нарушается
\end{answer}

\begin{answer}{6}{Даже если домохозяйства с конечным временным горизонтом сталкиваются со снижением налогов при своей жизни, но ожидают повышения налогов только после смерти, это не означает, что рикардианская эквивалентность однозначно нарушается.}
\qda. При наличии эффективного мотива оставления наследства (не углового решения) родители, зная, что налоги вырастут уже после их смерти (то есть лягут на плечи их потомков), увеличивают предназначенное потомкам наследство ровно настолько, чтобы компенсировать будущий рост налоговой нагрузки детей. В результате нынешнее поколение фактически интернализирует будущее налоговое бремя через альтруистическую связь поколений, и РЭ продолжает выполняться, несмотря на формально конечный горизонт жизни каждого отдельного домохозяйств
\end{answer}

\begin{answer}{7}{Наличие искажающего налогообложения может привести к нарушению рикардианской эквивалентности.}
\qda. Вывод РЭ был получен в предположении, что налоги аккордные (неискажающие), то есть их изменение не затрагивает предельные величины выбора агента. Если же вместо аккордных рассматривать искажающие налоги (например, пропорциональный налог на потребление или на трудовой доход), изменение их ставки во времени будет влиять на относительные (посленалоговые) цены благ в разные периоды и, следовательно, на условие первого порядка (правило Рамсея–Кейнса), меняя саму траекторию потребления, а не только уровень его финансирования. Значит, искажающее налогообложение - самостоятельный канал нарушения РЭ
\end{answer}

\begin{answer}{8}{Рикардианская эквивалентность справедлива для любого типа налогообложения. Являются ли налоги аккордными или искажающими не имеет значения.}
\qnet. Искажающие налоги (например, на потребление или на трудовой доход) затрагивают предельный выбор агента (правило Рамсея–Кейнса, решение о предложении труда), меняя приведённую стоимость доходов/расходов не нейтральным образом, приводя к нарушению нейтральности фискальной политики в отличие от аккордных налогов
\end{answer}

\begin{answer}{9}{Рикардианская эквивалентность может нарушаться, если часть домохозяйств сталкивается с ограничениями ликвидности.}
\qda. Базовая модель предполагает, что домохозяйства и правительство сберегают и занимают по одной и той же ставке $r$, то есть частный сектор имеет неограниченный доступ к заимствованию под ставку $r$. В реальности часть домохозяйств сталкивается с ограничением ликвидности: они не могут занять под ту же ставку (или вообще занять) даже при наличии высокой ожидаемой будущей доходности, и их текущее потребление ограничено текущим располагаемым доходом, а не приведённой стоимостью пожизненного дохода. Для таких домохозяйств временное снижение налогов сегодня увеличивает текущее потребление (поскольку снимает ограничение ликвидности здесь и сейчас), даже если приведённая стоимость их налоговых обязательств на протяжении жизни не изменилась, а значит, РЭ нарушается
\end{answer}

\begin{answer}{10}{По мнению П. Кругмана, в реальности США вряд ли имеют дело с ситуацией, когда действия спекулянтов на рынке государственных облигаций приводят к кризису государственного долга и ещё более глубокой рецессии.}
\qda. В своей книге \textit{«End This Depression Now!»} (2012) П. Кругман последовательно критикует концепцию так называемых \textit{«bond vigilantes»} («спекулянтов-бдителей» рынка облигаций) и доказывает, что сценарий подобного долгового кризиса в США крайне маловероятен.

\begin{enumerate}
    \item \textbf{Теоретическое обоснование (Институциональная специфика США):}
    Спекулянты на рынке облигаций — это инвесторы, сбрасывающие государственные ценные бумаги, когда они теряют уверенность в проводимой фискальной или монетарной политике, что приводит к скачку процентных ставок и росту стоимости заимствований. Однако Кругман указывает, что суверенное государство, которое:
    \begin{itemize}
        \item выпускает долг в собственной национальной валюте,
        \item обладает плавающим валютным курсом,
        \item имеет подконтрольный суверенный центральный банк (ФРС США), способный выступать кредитором последней инстанции,
    \end{itemize}
    в принципе не сталкивается с рисками неконтролируемого «набега» спекулянтов или вынужденного дефолта.

    \item \textbf{Эмпирический опыт Великой рецессии (2008–2012 гг.):}
    В период 2008–2011 гг. из-за кризиса США столкнулись с сочетанием низких налоговых поступлений и высоких государственных расходов, из-за чего федеральное правительство было вынуждено занимать рекордные суммы. 
    
    При каждом временном локальном росте ставок алармисты объявляли о «прибытии спекулянтов» и о том, что США больше не способны финансировать свой дефицит. Однако рыночная реальность опровергла эти страхи:
    \begin{itemize}
        \item Все временные всплески доходностей быстро сглаживались 
        \item Американские казначейские облигации ($US\ Treasuries$) сохраняли статус главного мирового «защитного актива» ($safe-haven\ asset$)
        \item К 2012 году доходности по заимствованиям США упали до исторических минимумов, несмотря на продолжающийся рост абсолютного размера госдолга
    \end{itemize}

    \item \textbf{Политико-экономический вывод:}
    Страх перед «спекулянтами-бдителями» Кругман считает необоснованным мифом, порождающим ложную «боязнь бюджетного дефицита». Опасения перед рынком облигаций отвлекают внимание от реальных экономических угроз — высокого уровня безработицы и падения совокупного спроса, а попытки сокращать дефицит во время рецессии лишь усугубляют спад
\end{enumerate}
\end{answer}

\begin{answer}{11}{П. Кругман объясняет колебания в доходности государственных облигаций США в 2009-2012 гг. изменениями в монетарной политике, регулирующей долгосрочную ставку процента.}
\qnet. В своей аргументации П. Кругман показывает, что колебания доходностей $US\ Treasuries$ определялись переоценкой рыночных ожиданий инвесторов относительно темпов восстановления экономики и выхода из ловушки ликвидности, а не прямым монетарным регулированием долгосрочных ставок со стороны ФРС.

\begin{enumerate}
    \item \textbf{Природа низкой доходности (Ловушка ликвидности):}
    Кругман объясняет низкий уровень доходностей казначейских облигаций слабым совокупным спросом, нахождением экономики в ловушке ликвидности ($liquidity\ trap$) и массовым «бегством в качество» ($flight\ to\ safety$) на фоне сокращения долговой нагрузки ($deleveraging$) частным сектором, а не целенаправленным манипулированием долгосрочной ставкой со стороны регулятора.

    \item \textbf{Механизм формирования долгосрочной ставки:}
    Федеральная резервная система напрямую контролирует лишь краткосрочную процентную ставку (которая в 2009–2012 гг. находилась на нулевом нижнем границе). Долгосрочные же ставки отражают средневзвешенные ожидания рынка относительно будущей траектории краткосрочных ставок на протяжении всего срока жизни облигации.

    \item \textbf{Причина колебаний в 2009–2012 гг.:}
    Колебания доходностей отражали смену настроений и прогнозов инвесторов относительно конца рецессии:
    \begin{itemize}
        \item Позитивные новости о восстановлении экономики порождали ожидания, что США скоро выйдут из ловушки ликвидности и ФРС начнет повышать краткосрочные ставки для предотвращения инфляции $\rightarrow$ инвесторы требовали большую компенсацию, и долгосрочные доходности росли
        \item Пессимистичные данные возвращали опасения стагнации $\rightarrow$ рынок закладывал сохранение нулевых краткосрочных ставок на длительный срок $\rightarrow$ долгосрочные доходности падали
    \end{itemize}
\end{enumerate}
\end{answer}

\begin{answer}{12}{П. Кругман согласен с экономистами, считающими, что в условиях ловушки ликвидности рост государственного долга приведёт к значительному росту ставки процента и помешает программе стимулирования со стороны ФРС США.}
\qnet. Позиция П. Кругмана прямо противоположна: он доказывает, что в условиях ловушки ликвидности рост государственного долга не приводит к повышению процентных ставок и не создает эффекта вытеснения, а является единственным эффективным способом восстановить совокупный спрос.

\begin{enumerate}
    \item \textbf{Природа ловушки ликвидности и избыток сбережений:}
    В условиях глубокого рецессионного шока частный сектор стремится сократить долговую нагрузку ($deleveraging$), из-за чего желаемые сбережения существенно превышают планируемые инвестиции ($S > I$). Даже при снижении краткосрочной процентной ставки до нулевой границы ($r \approx 0$) экономика остается в состоянии избытка праздно лежащих сбережений, которые не трансформируются в частные расходы

    \item \textbf{Отсутствие эффекта вытеснения ($Crowding-out\ effect$):}
    В нормальных экономических условиях государственный бюджетный дефицит может конкурировать с частным сектором за ограниченные заемные ресурсы, толкая процентные ставки вверх. Однако в ловушке ликвидности:
    \begin{itemize}
        \item Государственные заимствования лишь абсорбируют незадействованный частный капитал
        \item Бюджетный дефицит не вступает в конкуренцию с частными инвестициями
        \item Дополнительный выпуск облигаций не оказывает давления на ставки процента, оставляя их на рекордно низком уровне
    \end{itemize}

    \item \textbf{Согласованность с монетарным стимулированием:}
    Рост госдолга не только не мешает программе стимулирования ФРС США, но и дополняет её. Когда традиционные монетарные инструменты упираются в нулевую границу ставок, именно фискальная экспансия за счет привлечения неработающих сбережений позволяет вернуть совокупный спрос на потенциальный уровень
\end{enumerate}
\end{answer}

\begin{answer}{13}{По мнению П. Кругмана, после финансового кризиса 2008-2009 годов экономика США оказалась в ловушке ликвидности. В таких условиях увеличение государственного долга решает проблему избытка сбережений в частном секторе и может вывести экономику из рецессии.}
\qda. Данный тезис составляет центральную идею книги П. Кругмана \textit{«End This Depression Now!»}, где он доказывает необходимость активного фискального стимулирования в условиях балансового кризиса

\begin{enumerate}
    \item \textbf{Причина рецессии («Минский шок» и делеверидж):}
    Глубинная причина затяжного спада после кризиса 2008–2009 гг. заключалась в обвале кредитного пузыря. Обремененные долгами домохозяйства и корпорации были вынуждены одновременно начать процесс сокращения долговой нагрузки (\textit{deleveraging}) — резко урезать текущие расходы и наращивать сбережения, чтобы восстановить устойчивость своих балансов

    \item \textbf{Ловушка ликвидности и дисбаланс сбережений:}
    Массовое стремление частного сектора к сбережению создало сильный избыток предложения сбережений над спросом на них со стороны частных инвесторов ($S > I$). Поскольку ставка ФРС уже достигнула нулевой границы (\textit{Zero Lower Bound}, $r \approx 0$), экономика оказалась в ловушке ликвидности: рыночный механизм ценового регулирования через процентную ставку больше не мог автоматически уравновесить сбережения и инвестиции

    \item \textbf{Государство как «заемщик последней инстанции»:}
    В условиях, когда частный сектор категорически отказывается потреблять и инвестировать, единственный агент, способный восстановить совокупный спрос — это правительство. Наращивая государственный долг, государство:
    \begin{itemize}
        \item поглощает (абсорбирует) незадействованные избыточные сбережения частного сектора
        \item трансформирует праздно лежащие ресурсы в реальный совокупный спрос ($AD$) через госзакупки и инвестиционные программы
        \item замещает выпавшие расходы домохозяйств и бизнеса, предотвращая спираль дефляционного спада и выводя экономику из рецессии
    \end{itemize}
\end{enumerate}
\end{answer}

\begin{answer}{14}{П. Кругман считает, что так как рецессия 2008-2011 гг. обусловлена недостаточными сбережениями в частном секторе, выход из неё требует большей бережливости в государственном секторе, т.е. требует бюджетных профицитов.}
\qnet. П. Кругман показывает, что рецессия была обусловлена не дефицитом, а \textit{избытком} частных сбережений, а политика «бережливости» в государственном секторе (\textit{austerity}) лишь усугубляет экономический спад.

\begin{enumerate}
    \item \textbf{Избыток сбережений вместо их недостатка:}
    Причиной рецессии стал масштабный процесс сокращения долговой нагрузки (\textit{deleveraging}). На фоне высокого уровня докризисных долгов и неопределенности домохозяйства и бизнес одновременно начали сбрасывать обязательства, резко сокращать потребление и наращивать сбережения. В результате в экономике возник критический избыток сбережений и острая нехватка совокупного спроса

    \item \textbf{Опасность политики жесткой экономии (\textit{Austerity}):}
    Попытки правительства проявлять «бережливость» и стремиться к бюджетному профициту в условиях ловушки ликвидности Кругман называет катастрофической ошибкой. Сокращение государственных расходов в момент, когда частный сектор и без того урезает свои траты, ведет к дополнительному обвалу спроса, падению ВВП и росту безработицы

    \item \textbf{Парадокс бережливости (\textit{Paradox of Thrift}):}
    В рецессионных условиях срабатывает кейнсианский парадокс бережливости: когда абсолютно все агенты (и частный сектор, и государство) одновременно стремятся урезать расходы и увеличить сбережения, совокупный доход падает настолько сильно, что суммарный объем сбережений в экономике в итоге уменьшается, а депрессия только углубляется

    \item \textbf{Наращивание бюджетного дефицита:}
    Выходом из рецессии является не бюджетный профицит, а фискальная экспансия и рост государственного долга. Правительство должно изымать незадействованные избыточные сбережения с рынка через выпуск облигаций и направлять их на государственные закупки и социальную поддержку, замещая провал в частном спросе
\end{enumerate}
\end{answer}

\begin{answer}{15}{П. Кругман считает, что в условиях ловушки ликвидности сокращение государственных закупок будет воспринято как снижение налоговой нагрузки в будущем, что повысит доверие к политике и приведёт к росту ВВП.}
\qnet. Описанный тезис представляет собой концепцию «экспансионистской консолидации» (\textit{expansionary austerity}), которую П. Кругман жестко критикует и иронично называет мифом о «фее доверия» (\textit{confidence fairy}), доказывая её эмпирическую и теоретическую несостоятельность в условиях ловушки ликвидности.

\begin{enumerate}
    \item \textbf{Критика мифа о «фее доверия» (\textit{Confidence Fairy}):}
    Сторонники «экспансионистской строгой экономии» утверждают, что урезание бюджетных расходов сигнализирует рынку о снижении будущих налогов и повышает доверие бизнеса, что якобы стимулирует частные расходы и инвестиции. Кругман доказывает, что в условиях депрессии этот механизм не работает: фирмы и домохозяйства не увеличивают траты при отсутствии реального текущего спроса на продукцию, независимо от гипотетической будущей налоговой нагрузки

    \item \textbf{Разрушительный эффект в краткосрочном периоде ($SR$):}
    В ловушке ликвидности сокращение государственных закупок действует по прямому кейнсианскому каналу — оно пагубно урезает совокупный спрос ($AD$) и приводит к падению ВВП. При этом сократить дефицит бюджета не удается, так как сжатие экономики влечет за собой:
    \begin{itemize}
        \item резкое падение налоговых поступлений в бюджет
        \item рост обязательных расходов на программы социальной помощи (пособия по безработице)
        \item вероятное увеличение показателя долга к ВВП из-за опережающего падения знаменателя (ВВП)
    \end{itemize}

    \item \textbf{Урон долгосрочному потенциалу экономики ($LR$):}
    Политика жесткой экономии подрывает долгосрочные перспективы роста (проявляется эффект гистерезиса):
    \begin{itemize}
        \item из-за отсутствия спроса бизнес замораживает инвестиции и расширение производственных мощностей
        \item из-за циклической безработицы высококвалифицированные специалисты теряют навыки, дисквалифицируются и соглашаются на низкооплачиваемую работу или вовсе покидают рабочую силу
    \end{itemize}
\end{enumerate}
\end{answer}

\begin{answer}{16}{По мнению П. Кругмана, ограничительная фискальная политика в условиях рецессии может не улучшить, а ухудшить фискальное положение (в частности, отношение государственного долга к ВВП) даже в среднесрочной перспективе.}
\qda. Данный эффект П. Кругман описывает как концепцию «самопораженческой» политики жесткой экономии (\textit{self-defeating austerity}), доказывая, что урезание государственных расходов в разгар рецессии приносит результат, прямо противоположный декларируемому.

\begin{enumerate}
    \item \textbf{Структура показателя долговой нагрузки ($\frac{B}{Y}$):}
    Устойчивость государственных финансов определяется не абсолютным объемом долга $B$ (числитель), а соотношением долга и объема производимого ВВП — $\frac{B}{Y}$ (знаменатель). Динамика этого показателя зависит от того, какая из двух величин изменяется быстрее

    \item \textbf{Высокий мультипликатор и сжатие знаменателя ($Y$):}
    В условиях ловушки ликвидности ($r \approx 0$) центральный банк не может снизить процентные ставки для компенсации фискального сжатия. В этих обстоятельствах фискальный мультипликатор становится существенно больше единицы ($M > 1$). Сокращение государственных закупок ($G \downarrow$) вызывает мультипликативное падение совокупного спроса и резкое сжатие ВВП ($Y \downarrow\downarrow$)

    \item \textbf{Действие встроенных стабилизаторов и числитель ($B$):}
    Сжатие ВВП автоматически приводит к падению налоговых поступлений и одновременно увеличивает государственные трансферты (пособия по безработице и социальную помощь). В результате номинальный объем государственного долга $B$ сокращается значительно слабее, чем планировалось, либо продолжает расти из-за возникающего дефицита

    \item \textbf{Итоговый результат и эмпирические примеры:}
    Поскольку знаменатель ($Y$) падает быстрее и пропорционально сильнее, чем сжимается числитель ($B$), итоговое отношение $\frac{B}{Y}$ не снижается, а возрастает. В качестве эмпирического доказательства Кругман приводит опыт стран Южной Европы (Греция, Испания, Португалия) в 2010-х годах, где масштабное урезание госбюджетов привело к глубокой депрессии и историческим максимумам отношения госдолга к ВВП
\end{enumerate}
\end{answer}

\begin{answer}{17}{В основе аргументации П. Кругмана «Новые долги могут решить проблемы, созданные старыми долгами» лежит принцип рикардианской эквивалентности («новые и старые долги эквивалентны»).}
\qnet. Аргументация П. Кругмана прямо противоречит логике «чистой» рикардианской нейтральности долга и строится на принципиальной неоднородности экономических агентов, а не на предпосылке о репрезентативном индивиде.

\begin{enumerate}
    \item \textbf{Гетерогенность агентов и различия в склонности к потреблению ($MPC$):}
    Модель Кругмана базируется на разделении экономики на закредитованных должников (\textit{debtors}) и кредиторов (\textit{creditors}):
    \begin{itemize}
        \item Должники имеют крайне высокую предельную склонность к потреблению ($MPC \approx 1$) и вынуждены резко урезать расходы из-за необходимости гасить «старые» долги (\textit{deleveraging})
        \item Кредиторы имеют низкую предельную склонность к потреблению ($MPC \ll 1$) и склонны сберегать избыточные средства
    \end{itemize}
    Рикардианская эквивалентность, напротив, опирается на предпосылку о репрезентативном впередсмотрящем агенте, для которого структура и распределение долга внутри экономики не имеют значения

    \item \textbf{Механизм замещения долга (государственный вместо частного):}
    Тезис о том, что «новые долги могут решить проблемы старых», означает замещение частных обязательств государственными. Когда частный сектор сжимает расходы для выплаты старых долгов, совокупный спрос рушится. Государство, выпуская «новые» облигации, привлекает праздно лежащие сбережения кредиторов и направляет их на стимулирование спроса, предотвращая депрессию

    \item \textbf{Отрицание рикардианской нейтральности:}
    Если бы в экономике действовала строгая рикардианская эквивалентность, выпуск государственного долга был бы полностью нейтрален и не оказал бы никакого стимулирующего воздействия (население просто сберегало бы $100\%$ объема полученных средств в ожидании будущих налогов). Кругман категорически отрицает нейтральность госдолга, доказывая, что в условиях депрессии не все долги равноценны и правительственное заимствование эффективно компенсирует сжатие частного сектора
\end{enumerate}
\end{answer}

\begin{answer}{18}{П. Кругман считает, что списание плохих долгов в частном секторе и инфляция помогли бы экономике США быстрее выйти из рецессии.}
\qda. П. Кругман доказывает, что реструктуризация (списание) безнадежной задолженности и целевое повышение инфляционных ожиданий — это важнейшие механизмы, ускоряющие болезненный процесс делевериджа и стимулирующие совокупный спрос.

\begin{enumerate}
    \item \textbf{Списание плохих долгов частного сектора (\textit{Debt Relief}):}
    Рецессия обусловлена чрезмерным уровнем частного долга. Из-за различий в предельной склонности к потреблению ($MPC$) между закредитованными заемщиками ($MPC \approx 1$) и кредиторами ($MPC \ll 1$) сокращение расходов должниками обваливает совокупный спрос
    \begin{itemize}
        \item Списание части безнадежного долга снижает активы кредиторов, однако слабо влияет на их текущие расходы (так как кредиторы склонны сберегать)
        \item Для должников списание снижает долговую нагрузку и высвобождает доходы для текущего потребления
        \item В масштабах всей страны асимметричный эффект приводит к чистому росту совокупных расходов и ускоряет восстановление экономики
    \end{itemize}

    \item \textbf{Роль умеренной инфляции в ловушке ликвидности:}
    Когда номинальная процентная ставка упирается в нулевую границу ($i \approx 0$), традиционная монетарная политика теряет силу. В этих условиях инфляция создает два ключевых стимулирующих эффекта:
    \begin{itemize}
        \item \textbf{Отрицательная реальная процентная ставка:} Согласно эффекту Фишера, реальная ставка равна $r = i - \pi$. При $i = 0$ умеренно повышенная инфляция ($\pi > 0$) делает реальную ставку отрицательной ($r < 0$), что демотивирует удерживать деньги и подталкивает агентов к потреблению и инвестициям сегодня
        \item \textbf{Размывание фиксированного долга:} Инфляция обесценивает реальный размер накопленных номинальных долгов относительно будущих доходов, облегчая процесс восстановления балансов домохозяйств
    \end{itemize}

    \item \textbf{Итоговый вывод:}
    Сочетание частичного списания долгов и поддержания более высокой целевой инфляции позволяет сгладить негативные последствия балансового кризиса и выйти из депрессии значительно быстрее
\end{enumerate}
\end{answer}

\begin{answer}{19}{По мнению П. Кругмана, списание плохих долгов в частном секторе и инфляция только усугубляют рецессию. Лучшие средства это дефляция и полное погашение старых долгов.}
\qnet. Позиция П. Кругмана прямо противоположна: именно дефляция и попытки полного погашения долгов в разгар кризиса усугубляют депрессию, в то время как списание безнадежной задолженности и умеренная инфляция выступают ключевыми инструментами восстановления экономики.

\begin{enumerate}
    \item \textbf{Опасность дефляции и механизм «долговой дефляции» И. Фишера:}
    Требование полного погашения старых долгов в условиях спада запускает катастрофический механизм «долговой дефляции» (\textit{debt-deflation}). При дефляции ($\pi < 0$) реальная процентная ставка $r = i - \pi$ растет, а реальная стоимость зафиксированного в номинальном выражении долга увеличивается. Чем больше должники платят, тем тяжелее становится их реальный остаток долга, что заставляет их еще сильнее урезать расходы и загоняет экономику в спираль падения спроса

    \item \textbf{Смулирующий эффект списания плохих долгов:}
    Списание (реструктуризация) частных долгов напрямую увеличивает совокупные расходы за счет разницы в предельной склонности к потреблению ($MPC$):
    \begin{itemize}
        \item Списание уменьшает активы кредиторов и обязательства должников на одинаковую номинальную сумму
        \item Однако кредиторы имеют низкую $MPC$, поэтому падение их балансовых активов практически не влияет на их текущие потребление и расходы
        \item Должники обладают высокой $MPC \approx 1$, поэтому ослабление долгового пресса сразу же трансформируется в рост их реальных трат
    \end{itemize}
    В результате асимметрии реакций совокупный спрос в экономике увеличивается.

    \item \textbf{Роль умеренной инфляции:}
    В условиях ловушки ликвидности ($i \approx 0$) инфляция стимулирует экономику двумя путями:
    \begin{itemize}
        \item \textbf{Обеспечивает отрицательную реальную ставку ($r < 0$):} что подталкивает агентов от сбережения к тратaм и инвестициям
        \item \textbf{Размывает реальную стоимость долга:} инфляция «съедает» реальное бремя фиксированных номинальных обязательств, позволяя заемщикам быстрее восстановить платежеспособность без болезненного сжатия потребления.
    \end{itemize}
\end{enumerate}
\end{answer}

\begin{answer}{20}{По мнению П. Кругмана, самостоятельный выход экономики из рецессии затруднён тем, что «должники» не могут повысить расходы, а «кредиторы» не имеют стимулов тратить больше, чем раньше.}
\qda. П. Кругман доказывает, что асимметрия поведенческих стимулов между закредитованными заемщиками и кредиторами блокирует рыночные механизмы саморегулирования экономики и приводит к затяжному провалу совокупного спроса.

\begin{enumerate}
    \item \textbf{Вынужденное сжатие расходов должниками (\textit{Deleveraging}):}
    Домохозяйства и фирмы, накопившие высокий уровень задолженности в период докризисного бума, после обвала рынка активов столкнулись с угрозой неплатежеспособности. Обладая высокой предельной склонностью к потреблению ($MPC \approx 1$), должники вынуждены резко сокращать текущие расходы и направлять все высвободившиеся ресурсы на погашение старых долгов и восстановление балансов

    \item \textbf{Пассивность кредиторов и различия в $MPC$:}
    Кредиторы (более состоятельные домохозяйства и финансовые институты), получая процентные выплаты и возврат основного долга, не имеют стимула увеличивать собственное потребление на всю величину поступивших средств. Имея существенно более низкую предельную склонность к потреблению ($MPC \ll 1$), они направляют полученный поток ликвидности в сбережения

    \item \textbf{Провал рыночного саморегулирования:}
    В результате возникает системный дисбаланс: сокращение расходов со стороны должников не компенсируется соответствующим ростом трат со стороны кредиторов. В обычных условиях этот избыток сбережений мог бы быть сбалансирован падением процентной ставки, но в условиях ловушки ликвидности ($i \approx 0$) данный механизм не работает

    \item \textbf{Заключение:}
    Поскольку частный сектор заблокирован процессом сокращения долговой нагрузки, экономика оказывается в состоянии затяжной депрессии, из которой она не способна выйти самостоятельно без государственного фискального стимулирования
\end{enumerate}
\end{answer}

\begin{answer}{21}{По мнению П. Кругмана, фискальное стимулирование за счёт наращивания госдолга может решить проблему депрессивной экономики, в которой ограниченные в ликвидности потребители не могут восстановить расходы до желаемого уровня.}
\qda. П. Кругман показывает, что при наличии ограничений ликвидности у закредитованных потребителей именно фискальная экспансия за счет роста государственного долга выступает единственным эффективным механизмом восстановления совокупного спроса.

\begin{enumerate}
    \item \textbf{Ограничения ликвидности у должников:}
    Домохозяйства и фирмы, обремененные долгами, сталкиваются с жесткими ограничениями по ликвидности: они не могут получить новые кредиты и вынуждены урезать текущие расходы до минимума, чтобы погашать старые обязательства (\textit{deleveraging}). Они физически не способны восстановить свои расходы до желаемого уровня

    \item \textbf{Пассивность кредиторов и ловушка ликвидности:}
    Кредиторы, получающие средства в счет уплаты долгов, из-за нулевых процентных ставок и неопределенности не имеют стимулов увеличивать собственное потребление. Происходит накапливание невостребованных избыточных сбережений, а совокупный спрос в экономике обваливается

    \item \textbf{Роль государства в замещении спроса:}
    В этих условиях только правительство способно восполнить образовавшийся провал в экономике:
    \begin{itemize}
        \item Наращивая государственный долг, государство аккумулирует праздно лежащие сбережения кредиторов
        \item Направляя эти средства на государственные расходы и закупки ($G$), правительство трансформирует сбережения в реальный совокупный спрос ($AD$);
        \item Фискальный импульс замещает выпавшие расходы частного сектора и поддерживает ВВП, давая должникам возможность завершить процесс восстановления балансов без впадания экономики в затяжную депрессию
    \end{itemize}
\end{enumerate}
\end{answer}

\newpage
\section{Часть 3. Модель неоклассического экономического роста Рамсея}

\begin{answer}{1}{В модели Рамсея равновесная норма сбережений тем выше, чем выше темп роста населения.}
\qda. \begin{enumerate}
    \item \textbf{Условия стационарного состояния:}
    В стационарном состоянии удельные показатели потребления и капитала постоянны ($\dot{k} = 0$, $\dot{c} = 0$):
    \begin{itemize}
        \item Из уравнения динамики капитала $\dot{k} = f(k) - c - (n + \delta)k = 0$ следует стационарный уровень потребления:
        \[
        c^* = f(k^*) - (n + \delta)k^*
        \]
        \item Из уравнения Эйлера (при $\dot{c} = 0$) следует модифицированное золотое правило:
        \[
        f'(k^*) = \rho + n + \delta
        \]
    \end{itemize}

    \item \textbf{Вывод равновесной нормы сбережений ($s^*$):}
    По определению, норма сбережений — это доля выпуска, не направленная на текущее потребление:
    \[
    s^* = \frac{f(k^*) - c^*}{f(k^*)} = \frac{(n + \delta)k^*}{f(k^*)}
    \]
    Умножив и разделив данное выражение на $f'(k^*)$ и подставив условие Эйлера $f'(k^*) = \rho + n + \delta$, получаем:
    \[
    s^* = \frac{n + \delta}{\rho + n + \delta} \cdot \frac{k^* f'(k^*)}{f(k^*)} = \frac{n + \delta}{\rho + n + \delta} \cdot \alpha(k^*)
    \]
    где $\alpha(k^*) = \frac{k^* f'(k^*)}{f(k^*)}$ — эластичность выпуска по капиталу (доля капитала в национальном доходе).

    Для производственной функции Кобба–Дугласа доля капитала является константой ($\alpha(k^*) = \alpha$), поэтому окончательно получаем:
    \[
    \frac{\partial s^*}{\partial n} = \alpha \cdot \frac{(\rho + n + \delta)\cdot 1 - (n + \delta)\cdot 1}{(\rho + n + \delta)^2} = \frac{\rho \cdot \alpha}{(\rho + n + \delta)^2} > 0
    \]
    Поскольку субъективная норма дисконтирования домохозяйств $\rho > 0$ и доля капитала в доходе $\alpha > 0$, то равновесная норма сбережений $s^*$ положительно зависит от темпа роста населения $n$

    \item \textbf{Экономическая интуиция:}
    Более высокий темп роста населения $n$ усиливает эффект «размывания капитала». Чтобы удерживать капиталовооруженность труда на стационарном уровне $k^*$ для растущего числа новых работников, экономика вынуждена направлять более высокую долю совокупного дохода $s^*$ на инвестиции.
\end{enumerate}
\end{answer}

\begin{answer}{2}{В модели Рамсея равновесная норма сбережений тем ниже, чем выше норма субъективных межвременных предпочтений (т.е. чем более домохозяйства нетерпеливы).}
\qda. Из формулы стационарной нормы сбережений $s^* = \frac{(n+\delta)\alpha(k^*)}{\rho + n + \delta}$ видно, что рост $\rho$ (повышение нетерпеливости) увеличивает знаменатель, приводя к монотонному падению норме сбережений $s^*$
\end{answer}

\begin{answer}{3}{Т.к. домохозяйства нетерпеливы (дисконтируют полезность будущего потребления), равновесная норма сбережений в модели Рамсея ниже, чем норма сбережений, обеспечивающая соблюдение золотого правила накопления капитала в модели Солоу.}
\qda. Нетерпеливость домохозяйств ($\rho > 0$) в модели Рамсея приводит к выбору более низкого стационарного уровня капиталовооруженности по сравнению с золотым правилом Солоу, что обусловливает и более низкую равновесную норму сбережений.

\begin{enumerate}
    \item \textbf{Золотое правило накопления капитала в модели Солоу ($GR$):}
    Золотое правило максимизирует стационарное потребление на душу населения $c^* = f(k) - (n + \delta)k$. Условие первого порядка даёт:
    \[
    f'(k_{GR}) = n + \delta
    \]
    Соответствующая норма сбережений по золотому правилу равна доле капитала в доходе: $s_{GR} = \frac{(n+\delta)k_{GR}}{f(k_{GR})} = \alpha(k_{GR})$.

    \item \textbf{Модифицированное золотое правило в модели Рамсея:}
    В модели Рамсея домохозяйства эндогенно оптимизируют межвременное потребление, дисконтируя будущую полезность с нормой субъективного нетерпения $\rho > 0$. Из уравнения Эйлера в стационарном состоянии ($\dot{c} = 0$) следует:
    \[
    f'(k^*) = \rho + n + \delta
    \]

    \item \textbf{Сравнение капиталовооруженности ($k^*$ и $k_{GR}$):}
    Поскольку домохозяйства нетерпеливы ($\rho > 0$), предельная производительность капитала в модели Рамсея выше, чем в Золотом правиле Солоу:
    \[
    f'(k^*) = \rho + n + \delta > n + \delta = f'(k_{GR})
    \]
    В силу строгой вогнутости производственной функции ($f'' < 0$) предельный продукт капитала строго убывает по $k$. Следовательно:
    \[
    k^* < k_{GR}
    \]
    \includegraphics[width=0.8\linewidth]{/Users/alexandrapodlesnyuk/Desktop/Снимок экрана — 2026-09-19 в 17.01.06.png}
    \item \textbf{Сравнение норм сбережений ($s^*$ и $s_{GR}$):}
    Равновесная норма сбережений в модели Рамсея равна:
    \[
    s^* = \frac{(n + \delta)k^*}{f(k^*)} = \frac{n + \delta}{\rho + n + \delta} \cdot \alpha(k^*)
    \]
    Для производственной функции Кобба–Дугласа $\alpha(k^*) = \alpha(k_{GR}) = \alpha$. Так как $\rho > 0$, коэффициенты соотносятся как:
    \[
    \frac{n + \delta}{\rho + n + \delta} < 1 \implies s^* < s_{GR}
    \]
    Если бы домохозяйства были абсолютно терпеливы ($\rho = 0$), равновесные стационарные состояния и нормы сбережений в обеих моделях совпали бы ($k^* = k_{GR}$ и $s^* = s_{GR}$).
\end{enumerate}
\end{answer}

\begin{answer}{4}{Т.к. описание процесса накопления капитала в моделях Солоу и Рамсея совпадает, равновесная норма сбережений в модели Рамсея совпадает с нормой сбережений, обеспечивающей соблюдение золотого правила накопления капитала в модели Солоу.}
\qnet. Хотя уравнение динамики капитала $\dot{k} = f(k) - c - (n+\delta)k$ одинаково, в модели Рамсея норма сбережений выбирается эндогенно из задачи оптимизации. Из-за наличия $\rho > 0$ оптимальный выбор агентов даёт $k^* < k_{GR}$ и $s^* < s_{GR}$
\end{answer}

\begin{answer}{5}{Стационарная норма сбережений в модели Рамсея отрицательно зависит от нормы субъективных межвременных предпочтений, что следует из уравнения Эйлера: чем выше $\rho$, тем меньше домохозяйства ценят будущее потребление и меньше сбережения.}
\qda. Рост нормы субъективного дисконтирования $\rho$ отражает увеличение нетерпения домохозяйств, что ведет к снижению стационарной капиталовооруженности $k^*$ и монотонному падению равновесной нормы сбережений $s^*$.

\begin{enumerate}
    \item \textbf{Уравнение Эйлера и стационарное состояние:}
    Динамика удельного потребления описывается уравнением Эйлера:
    \[
    \frac{\dot{c}}{c} = \frac{1}{\sigma} \left( f'(k) - \delta - \rho - n \right)
    \]
    В стационарном состоянии ($\dot{c} = 0$) предельный продукт капитала равен:
    \[
    f'(k^*) = \rho + n + \delta
    \]
    Чем выше норма нетерпения $\rho$, тем выше требуемая равновесная доходность капитала. В силу убывающей предельной производительности ($f'' < 0$) рост $\rho$ снижает стационарный запас капитала $k^*$.

    \item \textbf{Вывод стационарной нормы сбережений ($s^*$):}
    Из уравнения динамики капитала $\dot{k} = f(k) - c - (n + \delta)k = 0$ следует, что стационарный объем инвестиций равен $(n + \delta)k^*$. Норма сбережений определяется как отношение инвестиций к выпуску:
    \[
    s^* = \frac{f(k^*) - c^*}{f(k^*)} = \frac{(n + \delta)k^*}{f(k^*)}
    \]
    Умножив и разделив данное выражение на $f'(k^*) = \rho + n + \delta$, получаем:
    \[
    s^* = \frac{n + \delta}{\rho + n + \delta} \cdot \frac{k^* f'(k^*)}{f(k^*)} = \frac{n + \delta}{\rho + n + \delta} \cdot \alpha(k^*)
    \]
    где $\alpha(k^*) = \frac{k^* f'(k^*)}{f(k^*)}$ — эластичность выпуска по капиталу (доля капитала в национальном доходе).

    Для производственной функции Кобба–Дугласа доля капитала постоянна ($\alpha(k^*) = \alpha$). Частная производная нормы сбережений по $\rho$ равна:
    \[
    \frac{\partial s^*}{\partial \rho} = \alpha (n + \delta) \cdot \left( -\frac{1}{(\rho + n + \delta)^2} \right) = -\frac{(n + \delta)\alpha}{(\rho + n + \delta)^2} < 0
    \]
    Производная строго отрицательна при любых положительных параметрах.

    \item \textbf{Экономическая интуиция:}
    При увеличении $\rho$ домохозяйства сильнее дисконтируют будущее полезность и отдают предпочтение текущему потреблению. Падение сбережений снижает предложение капитала, повышая реальную процентную ставку $r^* = f'(k^*) - \delta$. В результате снижаются и стационарный выпуск $f(k^*)$, и потребление $c^*$, однако выпуск падает сильнее необходимых для поддержания $k^*$ стационарных инвестиций, что приводит к однозначному снижению нормы сбережений $s^*$
\end{enumerate}
\end{answer}

\begin{answer}{6}{Стационарная норма сбережений в Рамсее положительно зависит от $\rho$: чем выше $\rho$, тем ниже $k^*$, а снижение капиталовооружённости сильнее снижает потребление, чем выпуск.}
\qnet. Зависимость стационарной нормы сбережений $s^*$ от параметра субъективного дисконтирования $\rho$ является отрицательной ($\frac{\partial s^*}{\partial \rho} < 0$), а приведенная в утверждении логика изменения выпуска и потребления инвертирована.

\begin{enumerate}
    \item \textbf{Отрицательный характер зависимости $s^*(\rho)$:}
    Из уравнения Эйлера в стационарном состоянии следует модифицированное золотое правило: $f'(k^*) = \rho + n + \delta$. Рост нетерпения домохозяйств ($\rho \uparrow$) снижает целевую капиталовооруженность $k^*$. Выражение для стационарной нормы сбережений имеет вид:
    \[
    s^* = \frac{(n + \delta)k^*}{f(k^*)} = \frac{n + \delta}{\rho + n + \delta} \cdot \alpha(k^*)
    \]
    Отсюда видно, что рост $\rho$ монотонно снижает норму сбережений ($\frac{\partial s^*}{\partial \rho} < 0$), а не увеличивает её.

    \item \textbf{Сравнение чувствительности выпуска и потребления к изменению $k^*$:}
    Стационарный выпуск на душу населения равен $y^* = f(k^*)$, а стационарное потребление — $c^* = f(k^*) - (n + \delta)k^*$. Возьмем производные обоих показателей по капиталу $k^*$:
    \begin{itemize}
        \item Изменение выпуска: $\frac{dy^*}{dk^*} = f'(k^*) = \rho + n + \delta$;
        \item Изменение потребления: $\frac{dc^*}{dk^*} = f'(k^*) - (n + \delta) = \rho$.
    \end{itemize}
    Поскольку $n + \delta > 0$, предельное падение выпуска при снижении $k^*$ строго больше предельного падения потребления:
    \[
    \frac{dy^*}{dk^*} = \rho + n + \delta > \rho = \frac{dc^*}{dk^*}
    \]
    Таким образом, снижение капиталовооруженности $k^*$ бьет по выпуску $y^*$ \textbf{сильнее}, чем по потреблению $c^*$.

    \item \textbf{Итоговый механизм изменения нормы сбережений:}
    Стационарную норму сбережений можно представить как $s^* = 1 - \frac{c^*}{y^*}$. Так как при падении $k^*$ выпуск сокращается в большей степени, чем потребление, доля потребления в доходе $\frac{c^*}{y^*}$ увеличивается. В результате норма сбережений $s^*$ снижается
\end{enumerate}
\end{answer}

\begin{answer}{7}{Стационарный уровень капиталовооружённости, соответствующий модифицированному золотому правилу в модели Рамсея, ниже уровня, соответствующего простому золотому правилу в модели Солоу.}
\qda. Домохозяйства в модели Рамсея эндогенно выбирают норму сбережений с учетом субъективного нетерпения ($\rho > 0$), что приводит к выбору более низкого стационарного запаса капитала по сравнению со статическим максимумом потребления в модели Солоу.

\begin{enumerate}
    \item \textbf{Простое золотое правило Солоу ($k_{GR}$):}
    Золотое правило Солоу определяет уровень капиталовооруженности, максимизирующий стационарное потребление на душу населения $c^* = f(k) - (n + \delta)k$. Условие первого порядка даёт:
    \[
    \frac{dc^*}{dk} = f'(k_{GR}) - (n + \delta) = 0 \implies f'(k_{GR}) = n + \delta
    \]

    \item \textbf{Модифицированное золотое правило Рамсея ($k^*$):}
    В модели Рамсея сбережения эндогенны и определяются межвременной оптимизацией полезности домохозяйств. В стационарном состоянии ($\dot{c} = 0$) из уравнения Эйлера следует:
    \[
    f'(k^*) = \rho + n + \delta
    \]

    \item \textbf{Математическое сравнение ($k^* < k_{GR}$):}
    Сравнивая предельные продукты капитала при условии положительной нормы субъективного дисконтирования ($\rho > 0$):
    \[
    f'(k^*) = n + \delta + \rho > n + \delta = f'(k_{GR})
    \]
    Так как производственная функция строго вогнута ($f''(k) < 0$), функция предельной производительности $f'(k)$ является строго убывающей. Следовательно:
    \[
    f'(k^*) > f'(k_{GR}) \implies k^* < k_{GR}
    \]
    \includegraphics[width=0.8\linewidth]{/Users/alexandrapodlesnyuk/Desktop/Снимок экрана — 2026-09-19 в 17.01.06.png}

    \item \textbf{Экономическая интуиция:}
    В модели Солоу норма сбережений задана экзогенно, а золотое правило указывает на технологический максимум устойчивого потребления без учета межвременных предпочтений. В модели Рамсея потребители сопоставляют полезность текущего и будущего потребления: из-за нетерпеливости ($\rho > 0$) они ценят сегодняшнее потребление выше завтрашнего. Поэтому им экономически невыгодно сберегать объем ресурсов, необходимый для накопления капитала до уровня $k_{GR}$, и равновесный капитал останавливается на уровне $k^* < k_{GR}$. Если бы домохозяйства не дисконтировали будущее ($\rho = 0$), оба золотых правила совпали бы ($k^* = k_{GR}$)
\end{enumerate}
\end{answer}

\begin{answer}{8}{Стационарная ставка процента, определяемая модифицированным золотым правилом в Рамсее, ниже стационарной ставки процента, соответствующей простому золотому правилу в Солоу.}
\qnet. Стационарная реальная ставка процента в модели Рамсея строго выше ставки, соответствующей золотому правилу Солоу, из-за наличия положительной нормы субъективного дисконтирования ($\rho > 0$).

\begin{enumerate}
    \item \textbf{Связь процентной ставки и предельного продукта капитала:}
    В условиях совершенной конкуренции реальная чистая ставка процента равна предельной производительности капитала за вычетом нормы выбытия:
    \[
    r = f'(k) - \delta
    \]

    \item \textbf{Ставка процента по Золотому правилу Солоу ($r_{GR}$):}
    Золотое правило максимизирует стационарное потребление, что даёт условие $f'(k_{GR}) = n + \delta$. Подставляя это выражение в формулу чистой ставки процента, получаем:
    \[
    r_{GR} = f'(k_{GR}) - \delta = (n + \delta) - \delta = n
    \]
    Таким образом, по золотому правилу Солоу стационарная ставка процента равна темпу роста населения $n$.

    \item \textbf{Ставка процента по модифицированному золотому правилу Рамсея ($r^*$):}
    В модели Рамсея из уравнения Эйлера в стационарном состоянии следует условие $f'(k^*) = \rho + n + \delta$. Чистая ставка процента равна:
    \[
    r^* = f'(k^*) - \delta = (\rho + n + \delta) - \delta = n + \rho
    \]

    \item \textbf{Математическое сравнение и экономическая интуиция:}
    Поскольку домохозяйства нетерпеливы ($\rho > 0$), сравнение двух ставок даёт:
    \[
    r^* = n + \rho > n = r_{GR}
    \]
    Из-за дисконтирования будущего полезность сегодняшнего потребления выше, поэтому домохозяйства сберегают меньше капитала ($k^* < k_{GR}$). В силу убывающей предельной производительности капитала ($f'' < 0$) меньший капитал приводит к более высокому предельному продукту $f'(k^*)$, что выражается в более высокой стационарной ставке процента $r^*$
\end{enumerate}
\end{answer}

\begin{answer}{9}{Равновесие в децентрализованной экономике в модели Рамсея является Парето-эффективным.}
\qda. В модели Рамсея децентрализованное конкурентное равновесие полностью совпадает с решением задачи централизованного социального планировщика (командным оптимумом), что означает выполнение первой теоремы благосостояния и парето-эффективность равновесной траектории.

\begin{enumerate}
    \item \textbf{Задача социального планировщика (командный оптимум):}
    Социальный планировщик максимизирует общественное благосостояние, выбирая траекторию удельного потребления $c(t)$ при физическом ресурсном ограничении экономики:
    \[
    \max_{\{c(t)\}} \int_{0}^{\infty} e^{-(\rho - n)t} u(c(t)) \, dt \quad \text{при} \quad \dot{k}(t) = f(k(t)) - (n + \delta)k(t) - c(t)
    \]
    Применяя метод Понтрягина (через приведенный гамильтониан $H = u(c) + \lambda [f(k) - (n + \delta)k - c]$), из условий первого порядка получаем дифференциальное уравнение динамики потребления:
    \[
    \frac{\dot{c}}{c} = \frac{1}{\sigma(c)} \left( f'(k(t)) - \delta - \rho \right)
    \]
    где $\sigma(c) = -\frac{c u''(c)}{u'(c)}$ — коэффициент относительной несклонности к риску.

    \item \textbf{Децентрализованное конкурентное равновесие:}
    В рыночной экономике агенты принимают решения самостоятельно:
    \begin{itemize}
        \item \textbf{Фирмы} на конкурентных рынках факторов производства оплачивают труд и капитал по их предельным продуктам:
        \[
        r(t) = f'(k(t)) - \delta, \quad w(t) = f(k(t)) - k(t)f'(k(t))
        \]
        \item \textbf{Домохозяйства} оптимизируют полезность при динамическом бюджетном ограничении:
        \[
        \dot{a}(t) = w(t) + r(t)a(t) - c(t) - n a(t)
        \]
        В равновесии активы на душу населения равны капиталовооруженности ($a(t) = k(t)$). Подставляя равновесные цены $w(t)$ и $r(t)$ в бюджетное ограничение:
        \[
        \dot{k}(t) = \left[ f(k(t)) - k(t)f'(k(t)) \right] + \left[ f'(k(t)) - \delta \right]k(t) - c(t) - n k(t) = f(k(t)) - (n + \delta)k(t) - c(t)
        \]
        Бюджетное ограничение домохозяйства полностью совпадает с ресурсным ограничением планировщика.
    \end{itemize}

    \item \textbf{Выполнение первой теоремы благосостояния:}
    Максимизация полезности домохозяйством при полученном бюджетном ограничении даёт тождественное уравнение Эйлера:
    \[
    \frac{\dot{c}}{c} = \frac{1}{\sigma(c)} (r(t) - \rho) = \frac{1}{\sigma(c)} \left( f'(k(t)) - \delta - \rho \right)
    \]
    Поскольку системы дифференциальных уравнений для $\dot{c}(t)$ и $\dot{k}(t)$, а также граничные условия (начальный капитал $k(0) = k_0$ и условия трансверсальности) у планировщика и рыночной экономики абсолютно идентичны, децентрализованное равновесие приводит к парето-эффективному распределению ресурсов. В системе отсутствуют рыночные несовершенства, экстерналии и искажающие налоги
\end{enumerate}
\end{answer}

\begin{answer}{10}{Равновесие в децентрализованной экономике в модели Рамсея необязательно совпадает с решением задачи поиска командного оптимума.}
\qnet. При отсутствии экстерналий и искажающих налогов децентрализованное рыночное равновесие строго совпадает с командным оптимумом
\end{answer}

\begin{answer}{11}{В модели Рамсея госзакупки, финансируемые аккордными налогами, приводят к полному вытеснению потребления в стационарном состоянии.}
\qda. В модели Рамсея постоянное увеличение государственных закупок, финансируемых аккордными (неискажающими) налогами, не изменяет стационарную капиталовооруженность $k^*$ и приводит к ровно эквивалентному падению частного потребления ($\Delta c^* = -\Delta g$), что означает полное вытеснение.

\begin{enumerate}
    \item \textbf{Ресурсное ограничение экономики с госсектором:}
    Пусть государственные закупки на душу населения $g(t)$ полностью финансируются за счет аккордных налогов $\tau(t) = g(t)$. Динамическое бюджетное ограничение домохозяйства с учетом равновесных цен факторов ($r = f'(k) - \delta$, $w = f(k) - k f'(k)$) принимает вид:
    \[
    \dot{k}(t) = w(t) + r(t)k(t) - \tau(t) - c(t) - n k(t) = f(k(t)) - (n + \delta)k(t) - c(t) - g(t)
    \]
    Таким образом, локус $\dot{k} = 0$ определяется как:
    \[
    c = f(k) - (n + \delta)k - g
    \]
    При росте $g > 0$ графическая кривая локуса $\dot{k} = 0$ в фазовом пространстве $(k, c)$ сдвигается строго вниз на величину $\Delta g$.

    \item \textbf{Уравнение Эйлера и неизменность капиталовооруженности ($k^*$):}
    Поскольку аккордный налог не искажает относительные цены и доходность капитала, оптимизационная задача домохозяйства приводит к стандартному уравнению Эйлера:
    \[
    \frac{\dot{c}}{c} = \frac{1}{\sigma(c)} \left( f'(k(t)) - \delta - \rho - n \right)
    \]
    В стационарном состоянии ($\dot{c} = 0$) равновесная капиталовооруженность определяется исключительно технологией и предпочтениями:
    \[
    f'(k^*) = \rho + n + \delta
    \]
    Величина $k^*$ (а следовательно, и стационарный выпуск $y^* = f(k^*)$) абсолютно не зависит от уровня государственных закупок $g$. Вертикальная линия локуса $\dot{c} = 0$ остается неподвижной.

    \item \textbf{Эффект полного вытеснения частного потребления ($\Delta c^* = -\Delta g$):}
    Подставляя стационарный капитал $k^*$ в условие $\dot{k} = 0$, получаем выражение для стационарного уровня потребления:
    \[
    c^* = f(k^*) - (n + \delta)k^* - g
    \]
    Возьмем полную производную стационарного потребления по государственным закупкам $g$:
    \[
    \frac{dc^*}{dg} = \underbrace{\left[ f'(k^*) - (n + \delta) \right]}_{=\,\rho} \cdot \underbrace{\frac{dk^*}{dg}}_{=\,0} - 1 = -1 \implies \Delta c^* = -\Delta g
    \]
    Увеличение государственных закупок на $1$ денежную единицу снижает стационарное частное потребление ровно на $1$ денежную единицу.

    \item \textbf{Экономическая интуиция:}
    Введение или рост аккордного налога создает чистый отрицательный эффект богатства для домохозяйств (снижает приведенную стоимость их располагаемого дохода). Потребители мгновенно сокращают свое текущее и будущее потребление на величину государственного изъятия, не меняя объем сбережений и накопительного капитала в долгосрочном периоде. В результате капитал $k^*$ и выпуск $y^*$ остаются прежними, а госсектор полностью вытесняет частное потребление
\end{enumerate}
\end{answer}

\begin{answer}{12}{В модели Рамсея госзакупки, финансируемые аккордными налогами, приводят к полному вытеснению инвестиций в стационарном состоянии.}
\qnet. В стационарном состоянии модели Рамсея государственные закупки вытесняют исключительно частное потребление ($\Delta c^* = -\Delta g$), в то время как стационарный объем инвестиций остаётся абсолютно неизменным ($\Delta i^* = 0$).

\begin{enumerate}
    \item \textbf{Стационарный капитал $k^*$ и уравнение Эйлера:}
    В стационарном состоянии ($\dot{c} = 0$) из уравнения Эйлера следует модифицированное золотое правило:
    \[
    f'(k^*) = \rho + n + \delta
    \]
    Данное равенство показывает, что стационарная капиталовооруженность $k^*$ определяется исключительно технологическими параметрами производственной функции $f(\cdot)$, нормой выбытия капитала $\delta$, темпом роста населения $n$ и субъективной нормой дисконтирования домохозяйств $\rho$. Величина $k^*$ фундаментально не зависит от уровня государственных закупок $g$.

    \item \textbf{Объем стационарных инвестиций ($i^*$):}
    Из условия постоянства капитала $\dot{k} = 0$ следует, что объем валовых инвестиций на душу населения, необходимый для покрытия выбытия и поддержания капиталовооруженности растущего населения, равен:
    \[
    i^* = (n + \delta)k^*
    \]
    Поскольку капиталовооруженность $k^*$ неизменна при любом уровне $g$, стационарный объем инвестиций также остается константой:
    \[
    \frac{di^*}{dg} = (n + \delta)\frac{dk^*}{dg} = 0 \implies \Delta i^* = 0
    \]
    Таким образом, в долгосрочном периоде вытеснения инвестиций не происходит вовсе

    \item \textbf{Сравнение с кейнсианской логикой (IS-LM) и вытеснением потребления:}
    \begin{itemize}
        \item В кратковременных кейнсианских моделях (например, IS-LM) рост госзакупок поднимает процентную ставку, что приводит к частичному вытеснению частных инвестиций
        \item В неоклассической модели Рамсея ставка процента в стационарном состоянии жестко фиксирована предложением капитала и нетерпеливостью домохозяйств ($r^* = \rho + n$). Поэтому рост государственных закупок полностью (на $100\%$) вытесняет только частное потребление ($\Delta c^* = -\Delta g$), оставляя инвестиции $i^*$, капитал $k^*$ и выпуск $y^*$ на прежнем уровне
    \end{itemize}
\end{enumerate}
\end{answer}

\end{document}